\section{Theory}\label{sec:theory}
% Source: G4/PROOFS.md §1-§5, §7 (self-review + independent review, 26/10/2026). Full proofs: Appendix A.
% Numbering here -> PROOFS.md: Lemma 1 = Bổ đề 1; Thm 1 = Định lý 1(a)(b)(c); Prop 1 = Mệnh đề 3;
% Thm 2 = Định lý 2; Prop 2 = Mệnh đề 4; Prop 3 = Mệnh đề 5 + 6.

We state what holds when the TFM is exact, what survives when it is not, and what the default sampler of
\cref{alg:icg} is \emph{not} covered by. Proofs are in \cref{app:proofs}.

\begin{lemma}[Full conditionals]\label{lem:fullcond}
Under the exchangeable prior, \eqref{eq:lemma1} holds for every $i$ and every $y$ with $p(y_{-i}\mid X)>0$.
\end{lemma}

\begin{theorem}[Exactness]\label{thm:exact}
Assume (A0)--(A3) and $\tfm=\ppd$. (a) For instance-level weak labels, single-site Gibbs updates
$\pi(y_i=k\mid y_{-i})\propto\tfm(k\mid x_i,D_{-i})\lik_i(k)$ satisfy detailed balance with respect to $\pi$, and the
chain is irreducible and aperiodic on $\mathrm{supp}(\pi)$, hence converges geometrically. (b) For LLP, the MH chain
that swaps two differently labelled points of one bag, with acceptance ratio computed from two forward passes, is
reversible and irreducible on $\mathrm{supp}(\pi)$. (c) The inductive estimator \eqref{eq:predict} converges almost
surely to $\bar p(y_*\mid x_*)=\sum_y\pi(y)\,p(y_*\mid x_*,(X,y))$, and
$\lvert\bar p-p(y_*\mid x_*,X,\W)\rvert\le\TV\big(p(y\mid X,\W),\,p(y\mid x_*,X,\W)\big)$, which is zero when the prior
is \emph{separable} (the distribution of the features carries no information about $p(y\mid x)$); the transductive chain converges to
$p(y_*\mid x_*,X,\W)$ exactly.
\end{theorem}

\begin{proposition}[Bayes optimality under the prior]\label{prop:bayes}
For any strictly proper scoring rule, $p(y_*\mid x_*,X,\W)$ uniquely minimises the Bayes risk over all measurable
predictors of $(X,\W,x_*)$, the expectation being taken under the TFM's prior.
\end{proposition}

\Cref{thm:exact} and \cref{prop:bayes} say that, with an exact TFM, \ours{} computes the Bayes answer for weak labels
under the TFM's prior \claim{3}\claim{10}. This is optimality \emph{on average over the prior}; on real data, which need
not follow it, the evidence is empirical (\cref{sec:exp}). A real TFM is only approximately the PPD of any prior, and
its conditionals need not even be compatible with a joint distribution. The next result is the one that applies to
the sampler as run.

\begin{theorem}[Stability under an approximate TFM]\label{thm:stability}
Let $P$ be the kernel of the algorithm with the exact PPD (stationary distribution $\pi$) and $\tilde P$ the same algorithm
with a TFM $\tfm>0$. Then $\tilde P$ has a unique stationary distribution $\tilde\pi$, even when the conditionals of $\tfm$
are incompatible, and
\begin{equation}\label{eq:stability}
  \TV(\tilde\pi,\pi)\le\tau(P)\,\bar\varepsilon,\qquad
  \tau(P)=\sum_{t\ge0}\delta(P^t),\quad \bar\varepsilon=\sum_y\tilde\pi(y)\,\TV\big(\tilde P(y,\cdot),P(y,\cdot)\big),
\end{equation}
where $\delta$ is Dobrushin's ergodicity coefficient. For the random-scan single-site kernel,
$\bar\varepsilon\le\min\{1,2\kappa\,\varepsilon_{\TV}\}$, where $\varepsilon_{\TV}$ bounds the TV error of the TFM's conditionals and $\kappa$ depends only on $\lik$; for
predictions, $\lvert\E_{\tilde\pi}\tfm(y_*\mid x_*,\cdot)-\bar p(y_*\mid x_*)\rvert\le\varepsilon'_{\TV}+\TV(\tilde\pi,\pi)$.
\end{theorem}

The form of \eqref{eq:stability} is classical perturbation theory for Markov chains
\citep{schweitzer1968,meyer1980,seneta1988,cho2001}, and it is tight in general (equality on a two-state chain). What
is specific to our setting is that the chain built from an incompatible TFM still has a unique stationary
distribution (it shares the transition graph of the exact chain), and the error terms are expressed through
quantities a TFM has, the TV error of its predictive. The bound reads as (integrated mixing time of the exact chain)
$\times$ (error per step): a TFM error matters most where the exact chain mixes slowly \claim{4}. Unlike the analysis of
iterative imputation by \citet{liu2014}, it is non-asymptotic and covers MH steps.

\begin{proposition}[Blocks]\label{prop:blocks}
(a) A block update whose proposal is the product of block-wise TFM predictives (for LLP, conditioned on bag counts),
accepted with the MH ratio computed by Brook's lemma along the changed labels, leaves $\pi$ invariant when $\tfm=\ppd$.
(b) Without the MH step, \cref{thm:stability} applies with an error $\bar\varepsilon_B$ that is bounded only under an
extra stability assumption on the PPD, and the bound grows with $n$ for fixed $K$.
\end{proposition}

\paragraph{Scope.} The configuration used in our experiments, uncorrected blocks with $K=10$ and bag-blocks for LLP, is
\emph{not} covered by a non-vacuous guarantee: \cref{thm:exact} and \cref{prop:bayes} apply to the single-site, swap and
MH-corrected samplers, and the bound of \cref{prop:blocks}(b) is vacuous at our sizes. We therefore report the
MH-corrected sampler as an ablation (\cref{sec:ablation}): at the same number of sweeps it does not predict better than
the default, and it is several times slower.

\paragraph{When can the sampler help?} The naive use of a TFM conditions it on one label per point drawn from the weak
labels. Consider the naive target $\tilde\eta_k(x)=P(\tilde w=k\mid x)$, the function a consistent learner fitted on
these naive labels approaches.

\begin{proposition}[When the naive TFM suffices]\label{prop:naive}
(a) If the naive label depends on $x$ only through $y$, with a channel whose off-diagonal entries are equal within each
column, then $\tilde\eta_k=b_k+\lambda_k\eta_k$ with $\lambda_k>0$: the naive target ranks points exactly as the true
posterior $\eta_k(x)=p(y=k\mid x)$, and in the binary case, with class prior $\varpi$, $\mathrm{AUC}(\eta)-\mathrm{AUC}(s)\le\E\lvert s-\tilde\eta\rvert/(\lambda\varpi(1-\varpi))$
for any score $s$. This covers PU-SCAR, class-dependent noise, missing-completely-at-random SSL, PLL with independent
distractors and LLP with random bags. (b) If the naive labels are obtained from the observation by a kernel that does not
depend on the parameter, their Fisher information is at most that of the observation, with equality iff the score of
the observation is a function of the naive labels. Equality holds for PU-SCAR and label noise; for PLL and LLP the
naive labels coarsen the observation and the loss is large (for LLP with bags of 8, about $3\%$ of the information in
the bag counts in a logistic model; \cref{app:proofs}).
\end{proposition}

\Cref{prop:naive} sorts weak labels into those where the naive TFM already uses all the information in the
observation, so that there is nothing for the sampler to recover through $\lik$ (PU-SCAR, class-dependent noise,
SSL), and those where it discards information (PLL, LLP). It predicted, before evaluation, where \ours{} could
outperform the naive TFM; analyses registered as exploratory in Study~1 agreed with this pattern \claim{12}
(\cref{app:secondary}).
